\subsection{Market structure in the core model \label{sec:MarketinCore}}
%putting CES and monopcomp in the context of market structure, and why it became popular

This section examines the background to the 'engine' of the \gls{coremodel}, built upon an earlier model of \gls{monocomp} first published by \citeauthor{dixit_monopolistic_1977} in 1977. In analytic approaches to market structure, there is a tight symbiotic relationship between theory structure and the mathematics that describe it. The particular form of market structure is very much determined by mathematical approach, and the \gls{coremodel} is no different. A full survey of the literature on market structure is beyond the scope of this thesis, but it is important to understand the context that led to the creation of the \gls{coremodel}. As well as shedding light on the model itself, arguments about market structure in analytic economics set the scene for discussing agent-based approaches to the same problems. 

Market structures in analytic models cannot deviate from the tracks laid down by the mathematical choices available. \cmt{A point here I make three times.} It is to the credit of Krugman especially that \gls{ge} has very consciously discussed these constraints and their implications. The argument this thesis makes is that \gls{abm} faces very different types of constraint; because of this, there is no clear way to 'translate' one into the other. 

Perfect competition and monopoly have straightforward mathematical interpretations: in the former, the size of the market is so large that no one economic actor can affect price - so prices are assumed to be fixed. Not only does this mean consumers can have no effect on price, but firms cannot either - and so strategic interaction is ruled out. In the latter, a monopoly firm can set a price that gives them the best revenue, based on the supply-demand curve they face for their product. Again, there are no other firms, so they face no strategic considerations. Between these two poles, a sprawling tangle of approaches to imperfect competition has been developed over the twentieth century; some of these are discussed below. But the impenetrability of this tangle of approaches in part explains why the publication of the \gls{dxs} is such a pivotal moment in this history: it provided a tractable method able to do away with many of the complexities facing theorists of imperfect competition. In particular, it could account for economies of scale, and it found a way around problems of strategic interaction.

Fujita, in his recent short history of spatial economics, makes some strong claims for the \gls{dxs}: von Thunen could not have produced anything other than the partial theory he did, any more than one could expect ``de Vinci, upon one's first view of his amazing `flying machines' drawn in the fifteenth century, to invent a real airplane'' without access to the internal combustion engine. For Fujita, the \gls{dxs} ``provided location theory with a powerful combustion engine.'' (\citealt{fujita_evolution_2010} p.8) That engine drives the \gls{coremodel} of \gls{ge}.




\subsection{Imperfect competition \label{sec:imperfect}}

problems related to dealing with imperfect competition have loomed over modern location theory. It is possible to build a general equilibrium model with perfect competition and linear production functions - but not if location is included. [why not?]







%%%%%%%%%%%%%%%
% CUTTINZ
%%%%%%%%%%%%%%%
\begin{comment}
It can be argued that many of the key distinctions in \gls{ge} flow from the maths: the clear blue water between perfect and imperfect competition models is a consequence of the maths, not of the real world. 


Now, in my margin notes, I've put ``vs `games you can play with CES functions' '' - presumably to sarcastically draw attention to the contrast here. An avowedly narrow model - general, in the sense of 'whole world', but still narrow. Is the engine up to the job?

It is worth examining how this was approached, and comparing to the use of another function that pre-dates the CES. The Cobb-douglas function, used earlier by x, can actually be derived as a special case of the CES. It also shares with the CES a first appearance in the literature where comparison to data was a strong motivation for its use. [find.]

---

\subsection{Market structure and space \label{Market structure and space}}

There are three broad categories of favoured market structure in analytic economics: perfect competition, monopoly and monopolistic competition. Why `favoured'? Each of these are analytically tractable. Perfect competition models can assume that actors all face the same prices, and can have no influence over them - they are too small compared to the extent of the market. Monopoly models can say how a single firm should set its prices to maximise revenue, when faced with a particular demand curve - and because there is only one firm, there are no strategic considerations. 

Outside of the two poles of perfect competition and monopoly, the \gls{monocomp} model first published by \citeauthor{dixit_monopolistic_1977}
in 1977 is by far the most popular approach to market structure in spatial settings. The approach itself is explained in-depth below, but before that it is important to grasp why the model was needed: it is impossible to combine perfect competition models with space. Why not? There are two related reasons. First: perfectly competitive models must work with linear output - no economies of scale (or `increasing returns') are allowed. To see why, consider a (rather artificial) situation where many firms have linear output, but one firm, producing the same good, is able to utilise economies of scale. They can thus produce each unit more cheaply than their competitors - and so this is no longer a perfect competition situation. A single price can no longer be assumed, nor can one say that no actor can influence price. In this case, that one firm immediately gains monopoly power. Second: it is impossible to include space in a model \emph{without} increasing returns - and so space and perfect competition models cannot be combined. Again, why not? Starrett (1978) is credited with first making this point formally. As Fujita and Thisse describe it -

\begin{quote}
``The essential question is whether the competitive price mechanism is able to explain the endogenous formation of economic agglomerations. To check the ability of a spatial model to do so, the best approach is to consider the case of a homogenous space in which economic agents are free to choose their locations... Starrett has shown that if space is homogenous and transport costly, then... the economy degenerates into separate single-location groups of agents with all trades taking place within, rather than between, groups. Consequently, the perfectly competitive price mechanism alone is unable to deal simultaneously with cities and trade.'' (\citealt{fujita_economics_2002} p.13-14) 
\end{quote}

To put it another way: if an actor is able to produce `in their own backyard', why would they choose to trade or collaborate? Any gain they would receive can never outweigh transport costs, if output in linear. However, with increasing returns introduced to the picture, the minimum conditions for spatial economic activity now exist. This `minimum condition for agglomeration' is investigated in an agent context through the 'backyard' series of models in this thesis. It is worth noting that, historically, 'backyard' production of various sorts has been a commonplace. Sheffield cutlers, for example, worked in this way, with many homes having their own mini-forges. [citation] 

In the discipline of economics, up until the publication of the Dixit/Stiglitz model, arguments about space in economic models had not moved on much from the 1940s and 1950s. As Krugman points out -

\begin{quote}
``... in the late 1980s mainsteam economists were almost literally oblivious to the fact that economies aren't dimensionless points in space.'' (Krugman, middle aged, p.1)
\end{quote}

This was at least in part because of earlier debates where economists had parted ways with a nascent regional science: the former insisted  that perfect competition models could include space implicitly as part of their cost structure; the latter argued that was nonsense. (Isard is one such, and it is interesting to note that he also seems to have little time for GE-style use of transport costs: ``... it is as unjustifiable and inadequate to lump transport costs into one category along with all trade resistances in the theory of trade, as it is to treat them as implicit production costs in the theory of production.'' (\citealt{isard_location_1956} p.478) After Krugman's reinterpretation of Dixit-Stiglitz in a spatial context, monopolistic competition gained two vital qualities appealling to mainstream economists: tractability and determinacy...

---
\end{comment}